\section{Introduction}\label{s:intro}

\newterm{Age-related macular degeneration} (AMD) represents a major clinical challenge, primarily affecting the \newterm{macula} of the retina and leading to a progressive decline in central vision \cite{coleman2008,lim2012}. The disease manifests initially through clinical indicators such as \newterm{drusen formation} and \newterm{retinal pigment epithelium irregularities}. As it advances, AMD is categorized into two forms: \newterm{neovascular} (wet or exudative) and \newterm{non-neovascular} (atrophic, dry, or non-exudative). Particularly in its late stages, AMD significantly impairs central visual acuity, frequently resulting in profound visual impairment and legal blindness. This deterioration in vision markedly affects individuals' quality of life and independence. With global prevalence projected to rise from approximately 200 million in 2020 to nearly 300 million by 2040, AMD presents a significant public health challenge with profound socioeconomic implications \cite{wong2014}.

% AMD is recognized as the third leading cause of severe and irreversible vision loss worldwide. However, the incidence of legal blindness and visual impairment due to AMD has seen a decline, attributed primarily to the advent of treatments that inhibit \newterm{\VEGF} (VEGF) \cite{lim2012}.\par

\newterm{\VEGF} (VEGF) plays a pivotal role in the pathogenesis of neovascular AMD. VEGF is a protein that promotes \newterm{angiogenesis}, the formation of new blood vessels, which is crucial for tissue repair and regeneration. However, in the context of AMD, overexpression of VEGF leads to the formation of abnormal blood vessels in the retina, contributing to the progression of the disease  \cite{farnoodian2017,witmer2003}. Consequently, the development of \newterm{\antiVEGF} (anti-VEGF) therapies, which involve the use of antibodies to inhibit the action of VEGF, has been a significant advancement in AMD treatment.

Recent research and therapeutic strategies in AMD have pivoted towards targeting angiogenesis, with anti-VEGF agents like Ranibizumab, Bevacizumab, and others emerging as cornerstones in this therapeutic domain. These agents work primarily by inhibiting the action of VEGF, consequently reducing the aberrant vascular growth associated with AMD. For instance, in the CATT study \cite{catt2012}, Ranibizumab and Bevacizumab demonstrated effectiveness in improving visual acuity over two years, highlighting the potential of anti-VEGF treatment \cite{catt2012}. Other studies, such as those by \citet{eldem2015} and \citet{ho2014}, also support the efficacy of Ranibizumab in improving visual outcomes for AMD patients. Additionally, emerging therapies and combinations, such as Ranibizumab with Indomethacin \cite{russo2018} and adjunctive therapies like \newterm{photodynamic therapy} \cite{semeraro2015}, offer promising avenues for enhanced treatment efficacy.

The primary goal of this systematic review is to meticulously appraise the efficacy, safety, and broader impacts of anti-angiogenic therapies in the context of AMD prevention, focusing on interventions like anti-VEGF agents and other innovative investigational approaches, with a particular emphasis on visual outcomes. Through the synthesis and critical assessment of the existant literature, this review aims to offer an exhaustive perspective on the current landscape of anti-angiogenic therapies for AMD.

\subsection{Research Objective}\label{s:researchobjective}
The objective of this research is to systematically review and critically evaluate the literature on anti-angiogenic therapies for AMD prevention, with a focus on assessing their efficacy, safety, and impact on visual outcomes, including different regimens, \eg \TreatAndExtend.

\subsection{Research Question}\label{s:researchquestion}
In patients with age-related macular degeneration, how do anti-angiogenic therapies influence visual acuity, retinal thickness, and the rate of neovascularization when compared to standard care or placebo?

\vfill
\newpage